\documentclass[11pt]{article}

\usepackage[margin=2.5cm]{geometry}
\usepackage{graphicx}
\usepackage{booktabs}
\usepackage{float}
\usepackage{caption}
\usepackage{hyperref}
\usepackage{biblatex}
\addbibresource{/home/thomas/bibliography/references.bib}

\title{First model  comparison}
\author{}
\date{\today}

\begin{document}

\maketitle

\section{Introduction}
This document presents the preliminary comparison of models for classifying music tracks by their genre. 
This first analysis has been run on the \texttt{FMA\_small} dataset\cite{fma_dataset}.
The main code used to run this analysis is \texttt{compare\_models.py}, along as \texttt{feature\_groups.py}. This code compare different classes of classifiers and rank their performances.

\section{The dataset}
The \texttt{FMA\_small} dataset is presented in \cite{fma_dataset}. It contains 8000 audio clips limited to 30 seconds (extracted from the middle of the tracks). Each audio clip is associated to a single top genre. The 8 different top genres are evenly distributed among the 8000 audio clips.

The dataset is split into a training set, which contains 6400 audio clips, a validation test and a test set which both contain 800 clips each. The genres are evenly distributed in the 3 subsets. The validation and the test set both are built in a way that their is no artist leakage in between the sets: each artist is bound to a given set and no artist is present over more than 1 set. This allows to train models on a given set and to test it on a different set with no overlapping artist. So the prediction is always done on unknown artists.

\section{The feature set}
blabla

\section{Model selection}
\subsection{The models}
The selection is done among different classes of models. The classes of models which have been tested are: logistic regression, K-nearest neighbors (KNN), support vector machine (SVM), random forest, extra trees, hist-gradient-boosting (HGB) and multilayer perceptron (MLP). All coming from the \texttt{sci-kit learn} library\cite{scikit-learn}.
For each class of models, a grid of parameters has been investigated. Here is a list of the parameters which have been explored:
\begin{itemize}
\item logistic regression: \texttt{C},
\item KNN: \texttt{n-neighbors}, \texttt{weights}, \texttt{p},
\item SVM: \texttt{C}, \texttt{gamma},
\item random forest: \texttt{max\_features}, \texttt{min\_samples\_leaf}, \texttt{n\_estimators}, \texttt{class\_weight},
\item extra trees: \texttt{max\_features}, \texttt{min\_samples\_leaf}, \texttt{n\_estimators}, \texttt{class\_weight},
\item HGB: \texttt{learning\_rate}, \texttt{max\_leaf\_nodes}, \texttt{min\_samples\_leaf}, \texttt{l2\_regularization},
\item MLP: \texttt{hidden\_layer\_sizes}, \texttt{alpha}.
\end{itemize}

\subsection{The training}
The training is done using cross-validation with the \texttt{StratifiedGroupKFold} method from the \texttt{sci-kit learn} package. This method creates stratified folds with non-overlapping groups. The non-overlapping group chosen here are the artists. In total, 5 folds were used to do the training of the models. It means that each model is trained on 4/5 of the data and tested on the remaining 1/5. This remaining 1/5 has no overlapping artist with the training set, so the prediction is done on unknown artists. Models are trained 5 times and at the end a prediction is done for all the data. This prediction is used to compare models.
The dataset used to do the cross-validation is the training set and the validation set combined.
The averaged F1 score over the 5 cross-validations is used to rank and select the models.


\section{Results}
Among all the models which have been test, 3 different classes of models were selected as promising. Here are the selected models with their parameters:
\begin{itemize}
\item SVM:
  \begin{itemize}
  \item   $\texttt{C}=1.3$, 
  \item $\texttt{gamma} = 0.000537$,
  \end{itemize}
\item HGB:
  \begin{itemize}
  \item   $\texttt{learning\_rate}=0.1$, 
  \item $\texttt{max\_leaf\_nodes}=31$, 
  \item $\texttt{min\_samples\_leaf}=10$, 
  \item $\texttt{l2\_regularization}=1$,
  \end{itemize}

\item MLP:
  \begin{itemize}
  \item  $\texttt{hidden\_layer\_sizes}=(64,32)$, 
  \item $\texttt{alpha}=10$.
  \end{itemize}
\end{itemize}

The table~\ref{tab:modelperf} sums up the performances of each model. It gives the F1 score measured on the 5 cross-validation folds (out-of-fold) and then the F1 score measured on the test set. There is a significant difference in the F1 score between the out of fold measurement and the one made on the test set.


\begin{table}[h]
  \centering
    \caption{The performances of each model. OOF stands for out-of-fold and TEST designates the performances measured on the test set.}
    \label{tab:modelperf}
  \begin{tabular}{lccc}
    \toprule
    & SVM & HGB & MLP \\
    \midrule
    F1\_score OOF & $0.554 \pm 0.015$  & $0.561 \pm 0.013$ & $0.551 \pm 0.017$ \\
    F1\_score TEST & 0.476 & 0.490 & 0.494 \\
    \bottomrule                        
  \end{tabular}
\end{table}


\begin{figure}[h]
    \centering
    \includegraphics[width=0.9\textwidth]{../../plotting/images/fma\_genre\_f1\_svm.png}
    \caption{The F1 scores for out of fold prediction and test prediction for the best SVM model.}
    \label{fig:svmgenres}
\end{figure}

\begin{figure}[h]
    \centering
    \includegraphics[width=0.9\textwidth]{../../plotting/images/fma\_genre\_f1\_hgb.png}
    \caption{The F1 scores for out of fold prediction and test prediction for the best HGB model.}
    \label{fig:hgbgenres}
\end{figure}

\begin{figure}[h]
    \centering
    \includegraphics[width=0.9\textwidth]{../../plotting/images/fma\_genre\_f1\_mlp.png}
    \caption{The F1 scores for out of fold prediction and test prediction for the best MLP model.}
    \label{fig:mlpgenres}
\end{figure}


This difference on the measured F1 score has been investigated. The figures~\ref{fig:svmgenres}, \ref{fig:hgbgenres} and~\ref{fig:mlpgenres} compare the F1 scores for the out-of-fold prediction and for the test prediction, for all the 8 genres. These figures show that the drop of performances is mainly due to a drop of performances to classify the Folk genre. The F1 score drops from 0.631 to 0.209 for SVM, from 0.627 to 0.199 for HGB and from 0.623 to 0.263 for MLP for this genre. There is also a significant drop of performance for the Experimental and Instrumental genres.

A model linked cause can be excluded since this behavior is observed for the 3 model classes.
So something specific to Folk (and partly Experimental and Instrumental) differs between the cross-validations and the test set.
Since the genres are evenly distributed on all the dataset, 
So the set used for the cross-validations (training + validation set) has the same genre distribution has the test set. So this difference in F1 scores is not coming from a difference on the genre distributions.







\printbibliography

\end{document}

%%% Local Variables:
%%% mode: LaTeX
%%% TeX-master: t
%%% End:
